\begin{theorem}[Residue-field formulas]\label{mod:delta-residueformula}
Let $F$ be a nonarchimedean local field, let $\psi$ have exact additive
conductor $n\in\mathbf Z$, and let $\gamma\in F^\times$ satisfy
\[
 \operatorname{ord}_F(\gamma)=n+1.
\]
Using the canonical Teichmuller section $[\,\cdot\,]_F:k_F\to\mathcal O_F$
from Definition~\ref{mod:localfield-residuefield}, define
\[
 \bar\psi_\gamma(x)=\psi([x]_F/\gamma)\qquad(x\in k_F).
\]
This is an additive character.  More precisely, for every $a\in\mathcal
O_F$ and every proof that $a$ is integral,
\[
 \bar\psi_\gamma(\bar a)=\psi(a/\gamma).
\]
Thus the value is independent of the integral lift, as in
Lemma~\ref{lem:residual-normalization}, and exactness of the additive
conductor proves $\bar\psi_\gamma\ne1$.

For exact local quasi-character data $\chi$, define a multiplicative
character $\bar\chi:k_F\to\mathbf C$ by evaluating $\chi$ on the canonical
Teichmuller lift of each element of $k_F^\times$ and extending by zero at
zero.  If $m_F(\chi)\leq1$, then for every $u\in\mathcal O_F^\times$,
\[
 \bar\chi(\bar u)=\chi(u).
\]
In particular, exact conductor one implies $\bar\chi\ne1$; no residue-field
unit representative is selected in either definition.

Let $G_\gamma(\chi,\psi)$ and
$\Delta_F^{\mathrm{fin}}(\chi,\psi;\gamma)$ be the representative-free
finite sum and computational local constant of
Definition~\ref{mod:delta-finitedefinition}.  If $m_F(\chi)=0$ and $\gamma$
is admissible, then
\[
 G_\gamma(\chi,\psi)=1,
 \qquad
 \Delta_F^{\mathrm{fin}}(\chi,\psi;\gamma)=\chi(\gamma).
\]
This is the exact conductor-zero formula: the quotient $U_F^0/U_F^0$ is a
singleton, the additive value of its sole summand is one, and the unit value
of $\chi$ is one.

Suppose instead that $m_F(\chi)=1$ and $\gamma$ is admissible.  Then
$\operatorname{ord}_F(\gamma)=n+1$, so the preceding residual additive
character applies.  With the normalization of
Definition~\ref{mod:finitefield-characterapi},
\[
 \tau_{k_F}(\bar\chi,\bar\psi_\gamma)
 =-\sum_{x\in k_F^\times}
     \bar\chi(x)^{-1}\bar\psi_\gamma(x),
\]
the exact local-to-residue identities are
\[
 G_\gamma(\chi,\psi)
   =-\tau_{k_F}(\bar\chi,\bar\psi_\gamma),
\]
and
\[
 \Delta_F^{\mathrm{fin}}(\chi,\psi;\gamma)
 =-\chi(\gamma)\,
   \operatorname{ph}\!\left(
     \tau_{k_F}(\bar\chi,\bar\psi_\gamma)\right).
\]
The proof reindexes the local quotient by the canonical isomorphism
$U_F^0/U_F^1\simeq k_F^\times$, then uses the arbitrary-lift evaluation
theorems term by term.  Hence the displayed minus sign, multiplicative
character inversion, admissible-denominator factor, and phase are all
literal consequences of the definitions; no desired finite-field formula
is assumed.
\LeanFile{LanglandsFirstMainLemma/Delta/ResidueFormula.lean}
\lean{LanglandsFirstMainLemma.residualAddChar}
\lean{LanglandsFirstMainLemma.residualAddChar_integral_lift}
\lean{LanglandsFirstMainLemma.residualAddChar_residueUnits}
\lean{LanglandsFirstMainLemma.residualAddChar_ne_one}
\lean{LanglandsFirstMainLemma.teichmullerLocalUnits}
\lean{LanglandsFirstMainLemma.residueUnits_teichmullerLocalUnits}
\lean{LanglandsFirstMainLemma.residualMulChar}
\lean{LanglandsFirstMainLemma.residualMulChar_apply_unit}
\lean{LanglandsFirstMainLemma.residualMulChar_residueUnits}
\lean{LanglandsFirstMainLemma.residualMulChar_ne_one}
\lean{LanglandsFirstMainLemma.finiteGaussSum_conductor_zero}
\lean{LanglandsFirstMainLemma.deltaFinite_conductor_zero}
\lean{LanglandsFirstMainLemma.finiteGaussSum_conductor_one}
\lean{LanglandsFirstMainLemma.deltaFinite_conductor_one}
\leanok
\SourceText{Lemma lem:residual-normalization; Proposition prop:tame-conductor-one-comparison; Proposition prop:unramified-new.}
\uses{mod:delta-finitedefinition,mod:delta-elementary,mod:finitefield-characterapi,mod:localfield-residuefield}
\FormalizationContract The residual characters use the canonical Teichmuller
section only for their definitions; separate theorems identify their values
with every integral or unit lift.  The conductor-one finite sum is proved by
an explicit quotient equivalence and is exactly the negative of Langlands's
normalized Gauss sum.  The local-constant theorem then transports that minus
sign through the nonzero phase, retaining $\chi(\gamma)$ and the inverse on
$\bar\chi$ with no normalization or representative hypothesis suppressed.
\end{theorem}
